\itemhead{2.3 例题 · 例题8（教材p.56）}{XJ-B1-C02-EX-P056-Y1002}
\label{XJ-B1-C02-EX-P056-Y1002}
{\footnotesize 湘教版教材p.56 / PDF第63页 / 2.3 例题 / 例题8}
\field{原题}{已知汽车从踩刹车到停车所滑行的距离 $L$（$\mathrm m$）与速度 $v$（$\mathrm{km/h}$）之间有如下关系式：$L=k\cdot M\cdot v^2$，其中 $k$ 是比例系数，且 $k>0$，$M$ 是汽车质量（$\mathrm t$）。若一辆卡车不装货物（司机体重忽略不计）以 $36\,\mathrm{km/h}$ 的速度行驶时，从刹车到停车需要走 $20\,\mathrm m$。当这辆卡车装着等于车重的货物行驶时，为保证安全，要在发现前面 $20\,\mathrm m$ 处有障碍物时能在离障碍物 $5\,\mathrm m$ 以外处停车，则最高速度应低于多少（设司机发现障碍物到踩刹车需经过 $1\,\mathrm s$）？}
\field{参考答案}{载货后刹车距离为 $5v^2/162$，反应距离为 $5v/18$。安全条件 $5v^2/162+5v/18<15$，得 $0\le v<18$。最高速度应低于18 km/h；“以外”对应严格不等式。}
{\small \field{关联知识}{I04 一元二次不等式与参数分类（教材帮p.64、65、66 / PDF73、74、75）；F03 函数求值、解析式与函数建模（教材帮p.80、87、88、89、90、131 / PDF89、96、97、98、99、140）}
\field{方法与易错点}{化标准式、查系数、求根或判别式、按图定号，最后与实际范围取交集。 易错：套“内小外大”却忘开口；丢重根或降次；漏整数产量限制。}
\field{答案依据}{例题原解位置：湘教版教材p.56及题库原图；独立编写可编辑参考解答。答案册此部分主要收练习、习题，不伪造例题答案册位置。}}
